
\documentclass[12pt,a4paper]{article}

\usepackage[utf8]{inputenc}
\usepackage[T1]{fontenc}
\usepackage{amssymb}
\usepackage[brazil]{babel}
\usepackage{lmodern}
\usepackage{geometry}
\usepackage{setspace}
\usepackage{parskip}
\usepackage{enumitem}
\usepackage{booktabs}
\usepackage{longtable}
\usepackage{xcolor}
\usepackage{listings}
\usepackage{hyperref}
\usepackage{amsmath}
\usepackage{graphicx}
\usepackage{float}

\geometry{margin=2.5cm}
\onehalfspacing
\hypersetup{colorlinks=true,linkcolor=blue,urlcolor=blue,citecolor=blue}

\definecolor{codegray}{gray}{0.95}
\lstset{
  backgroundcolor=\color{codegray},
  basicstyle=\ttfamily\small,
  breaklines=true,
  frame=single,
  columns=fullflexible,
  keepspaces=true
}

\title{\textbf{Guia Completo de Execução do PIBITI}\\
\large Detecção de deepfakes de áudio em português brasileiro por decomposição Wavelet Packet e descritores estatísticos do contorno de onda das sub-bandas}

\author{Thiago Mateus Schiavi\\Fatec Ourinhos -- ADS}
\date{2026--2027}

\begin{document}
\maketitle
\tableofcontents
\newpage

\section{Visão de 5 minutos}

Este manual transforma o Plano de Atividades oficial do PIBITI em um roteiro prático de execução. O projeto tem como núcleo a engenharia de \textit{features}: aplicar Wavelet Packet Transform (WPT), reconstruir sub-bandas no domínio do tempo e extrair descritores estatísticos do contorno de onda.

O núcleo mínimo definido no plano é:
\begin{itemize}
    \item nível WPT $L=4$;
    \item wavelet-mãe db4;
    \item dois corpora, sendo um de maior porte e BraSpoof-50;
    \item classificadores GMM e XGBoost;
    \item avaliação por EER e min-DCF;
    \item execução em CPU.
\end{itemize}

O fluxo geral é:

\begin{center}
Áudio $\rightarrow$ Partição $\rightarrow$ Padronização $\rightarrow$ WPT
$\rightarrow$ Sub-bandas $\rightarrow$ Reconstrução $\rightarrow$ Features
$\rightarrow$ Vetor de Features $\rightarrow$ Classificador
$\rightarrow$ Score $\rightarrow$ EER/min-DCF $\rightarrow$ Análise
\end{center}

O espaço completo $L\in\{4,5\}\times$ 3 wavelets-mãe $\times$ 4 corpora $\times$ 4 classificadores é uma meta de expansão. A entrega mínima é o núcleo descrito acima.

\section{Como usar este guia}

Há três categorias de informação:
\begin{description}
    \item[Definido no plano] aquilo que aparece como atividade, marco ou orientação no documento oficial.
    \item[Sugestão prática] organização adicional proposta neste guia para facilitar a execução.
    \item[Conhecimento complementar] explicações técnicas necessárias para compreender ou implementar o projeto.
\end{description}

Quando o plano não define uma divisão diária ou semanal detalhada, este guia não deve apresentar uma sugestão como se fosse uma exigência oficial.

\section{Cronograma oficial dos 12 meses}

\begin{longtable}{p{1.2cm}p{2.2cm}p{6cm}p{4cm}}
\toprule
Mês & Período & Foco principal & Marco/entrega\\
\midrule
1 & set/2026 & Fundamentação (wavelet/WPT) + ambientação & Ambiente + repositório Git\\
2 & out/2026 & Partição + teoria WPT & Partição validada + fichamento\\
3 & nov/2026 & Revisão + WPT-piloto (reconstrução) & Decomposição WPT funcionando\\
4 & dez/2026 & Calibração de nível/wavelet + início dos descritores & Sub-bandas + primeiros descritores\\
5 & jan/2027 & Extração completa dos descritores & Vetor de features do núcleo\\
6 & fev/2027 & Treino do núcleo + biblioteca v0 + relatório parcial & Relatório parcial + núcleo\\
7 & mar/2027 & Expansão SVM/RF, db8/sym8, L=5 & Matriz ampliada\\
8 & abr/2027 & Demais corpora + custo em CPU & Cobertura ampliada + custo\\
9 & mai/2027 & Avaliação completa + contribuição de sub-bandas & Resultados + análise\\
10 & jun/2027 & Biblioteca consolidada + repositório + redação & Biblioteca documentada + rascunho\\
11 & jul/2027 & Submissão + finalização do artefato & Artigo submetido + repositório\\
12 & ago/2027 & Relatório final + encerramento & Relatório final + apresentação\\
\bottomrule
\end{longtable}

\section{Mês 1 -- Fundamentação e ambientação}

\subsection{Semana 1 -- Panorama do problema}

O plano determina o estudo de bona fide versus spoof, protocolos ASVspoof, métricas EER/min-DCF e a criação do repositório Git.

\textbf{Checklist:}
\begin{itemize}[label=$\square$]
\item Entender bona fide e spoof.
\item Entender a finalidade de EER e min-DCF.
\item Criar o repositório.
\item Criar data/, splits/, src/, features/, results/ e reports/.
\item Registrar as primeiras referências.
\end{itemize}

\subsection{Semana 2 -- Wavelets e Wavelet Packet}

\textbf{Esta é a semana atual.}

O plano determina estudar a árvore de decomposição, nível $L$, wavelet-mãe e o motivo pelo qual a WPT refina todas as sub-bandas.

\textbf{Estudar:}
\begin{itemize}
\item conceito de wavelet;
\item wavelet-mãe;
\item decomposição;
\item sub-bandas;
\item nível $L$;
\item DWT versus WPT;
\item árvore WPT.
\end{itemize}

Para WPT, uma decomposição de nível $L$ produz $2^L$ sub-bandas no nível final. Assim, para $L=4$, o projeto trabalha com 16 sub-bandas.

\textbf{Exercício recomendado:} fazer uma implementação pequena com PyWavelets apenas para observar a árvore e seus nós. Nesta semana, não é necessário processar o corpus inteiro nem treinar classificadores.

\begin{lstlisting}[language=Python]
import pywt
import soundfile as sf

audio, sr = sf.read("seu_audio.wav")

wp = pywt.WaveletPacket(
    data=audio,
    wavelet="db4",
    mode="symmetric",
    maxlevel=4
)

nodes = wp.get_level(4, order="freq")

print("Quantidade de sub-bandas:", len(nodes))

for i, node in enumerate(nodes):
    print(i, node.path, len(node.data))
\end{lstlisting}

\textbf{Resultado esperado:} identificação dos 16 nós/sub-bandas no nível 4.

\subsection{Semana 3 -- Ambiente Python}

O plano prevê pywavelets, numpy, scipy, scikit-learn, xgboost, soundfile e librosa, com execução em CPU.

\begin{lstlisting}
python -m venv .venv
.venv\Scripts\activate
pip install pywavelets numpy scipy scikit-learn xgboost pandas matplotlib soundfile librosa
pip freeze > requirements.txt
\end{lstlisting}

\subsection{Semana 4 -- Fichamento e pergunta de pesquisa}

Iniciar o fichamento e registrar a pergunta de pesquisa prevista no plano: se descritores do contorno de onda de sub-bandas WPT associados a ML clássico conseguem detectar fala sintética PT-BR com bom desempenho e baixo custo.

\section{Mês 2 -- Partição e teoria WPT}

\textbf{Objetivo oficial:} garantir uma partição válida e dominar a teoria WPT.

A partição deve ser única e congelada. Deve-se verificar integridade, ausência de vazamento de locutor e balanceamento por corpus.

\textbf{Checklist:}
\begin{itemize}[label=$\square$]
\item Verificar/gerar a partição.
\item Confirmar 16 kHz e mono quando a padronização for responsabilidade do projeto.
\item Verificar labels.
\item Verificar speaker leakage.
\item Confirmar o subconjunto dimensionado para CPU.
\item Consolidar o fichamento WPT.
\end{itemize}

\section{Mês 3 -- WPT-piloto e reconstrução}

Implementar WPT configurável e reconstrução das sub-bandas no domínio do tempo com db4.

\textbf{Regra crítica:} validar primeiro com sinais sintéticos conhecidos. Verificar reconstrução quase perfeita, conservação de energia, alinhamento e comprimento. O plano destaca padding e bordas como pontos de atenção.

\section{Mês 4 -- Nível, wavelet e primeiros descritores}

Começar em $L=4$ e preparar a comparação futura entre db4, db8 e sym8. Implementar inicialmente energia/RMS e ZCR.

O envelope deve ser obtido pela transformada de Hilbert; seus descritores incluem média, desvio e curtose.

\section{Mês 5 -- Extração completa}

Implementar, por sub-banda:
\begin{itemize}
\item energia/RMS;
\item ZCR;
\item envelope: média, desvio e curtose;
\item entropia;
\item fator de forma;
\item fator de crista.
\end{itemize}

Concatenar os resultados das sub-bandas para produzir um vetor por áudio. Gerar a matriz de features do núcleo e executar sanity checks, incluindo ausência de NaN e reprodutibilidade por seed.

\section{Mês 6 -- Treino do núcleo}

Treinar GMM e XGBoost sobre os vetores do núcleo. Avaliar EER, min-DCF e custo em CPU.

Para min-DCF, o plano fixa:
\[
C_{\mathrm{miss}}=1,\qquad C_{\mathrm{fa}}=1,\qquad P_{\mathrm{tar}}=0.05.
\]

Priorizar EER em razão do desbalanceamento. Entregar relatório parcial, núcleo validado e biblioteca v0.

\section{Mês 7 -- Expansão}

Somente depois de o núcleo estar validado:
\begin{itemize}
\item adicionar SVM e Random Forest;
\item usar StandardScaler para SVM;
\item testar db8 e sym8;
\item testar $L=5$;
\item registrar corpus, nível, wavelet e classificador.
\end{itemize}

\section{Mês 8 -- Corpora adicionais e CPU}

Incorporar FakeBrAccent, Fake Voices e o que couber dentro do custo computacional. Medir tempo de decomposição, extração, treino, inferência e memória.

\section{Mês 9 -- Avaliação e contribuição das sub-bandas}

Calcular EER e min-DCF para as combinações previstas. Investigar contribuição de sub-bandas e descritores usando feature importance e permutation importance. Produzir curvas DET/ROC e gráficos comparativos.

\section{Mês 10 -- Biblioteca e artigo}

Consolidar a biblioteca WPT + descritores, documentar o repositório, organizar README, environment.yml e resultados. Não redistribuir áudios; manter código, protocolos e resultados agregados.

\section{Mês 11 -- Submissão}

Revisar e submeter o artigo. Congelar biblioteca e repositório, adicionar LICENSE e CITATION.cff quando aplicável e verificar licenças dos corpora/modelos.

\section{Mês 12 -- Relatório final}

Redigir objetivos alcançados, metodologia, resultados por corpus/nível/wavelet, análise de sub-bandas, custo em CPU, limitações e trabalhos futuros. Preparar apresentação/pôster.

\section{Fundamentos de áudio}

O pipeline deve padronizar os dados conforme o escopo do projeto. O plano cita 16 kHz e mono na frente de partição/padronização.

Exemplo de verificação:

\begin{lstlisting}[language=Python]
import soundfile as sf

audio, sr = sf.read("audio.wav")

print("Sample rate:", sr)
print("Shape:", audio.shape)
print("Canais:", 1 if audio.ndim == 1 else audio.shape[1])
print("Amostras:", len(audio))
print("Duracao:", len(audio) / sr, "segundos")
\end{lstlisting}

\section{WPT em detalhes}

A WPT organiza a decomposição em uma árvore. Em nível 4, há 16 sub-bandas finais:

\[
2^4=16.
\]

A utilização de WPT permite refinar todas as sub-bandas, em contraste com a decomposição convencional que não mantém essa mesma expansão completa.

A implementação inicial deve usar db4 e permitir que o nível seja configurável.

\section{Reconstrução}

Cada sub-banda deve ser reconstruída no domínio do tempo. Antes de aplicar o método a grandes corpora, testar:
\begin{enumerate}
\item sinal sintético;
\item decomposição;
\item reconstrução;
\item comparação com o sinal original;
\item erro de reconstrução;
\item comprimento;
\item energia;
\item comportamento nas bordas.
\end{enumerate}

\section{Features}

\subsection{RMS e energia}

Para uma sequência $x[n]$ com $N$ amostras:
\[
RMS=\sqrt{\frac{1}{N}\sum_{n=1}^{N}x[n]^2}.
\]

A energia pode ser calculada como:
\[
E=\sum_{n=1}^{N}x[n]^2.
\]

\subsection{ZCR}

A taxa de cruzamentos por zero mede a frequência com que o sinal muda de sinal.

\subsection{Envelope de Hilbert}

O sinal analítico pode ser obtido por Hilbert. O envelope é a magnitude desse sinal analítico. O plano determina extrair média, desvio e curtose do envelope.

\subsection{Entropia}

A entropia é usada como descritor e aparece também no contexto de seleção de base WPT.

\subsection{Fator de forma}

O fator de forma é relacionado à razão entre RMS e média retificada, conforme a definição adotada e documentada na implementação.

\subsection{Fator de crista}

\[
CF=\frac{\max |x[n]|}{RMS}.
\]

\section{Vetor de features}

O princípio é:

\[
1\ \text{áudio} \rightarrow 16\ \text{sub-bandas}
\rightarrow \text{descritores} \rightarrow 1\ \text{vetor}.
\]

A quantidade exata de colunas deve ser determinada a partir dos descritores efetivamente implementados e documentada. Não se deve assumir uma dimensão diferente sem registrar a decisão.

\section{Machine Learning}

\subsection{GMM}

O núcleo prevê dois modelos, bona fide e spoof, e uma razão de verossimilhança para produzir score.

\subsection{XGBoost}

XGBoost será usado como classificador do núcleo. Seus hiperparâmetros e seed devem ser registrados.

\subsection{Expansão}

SVM e Random Forest pertencem à expansão do mês 7. A padronização é especialmente importante para SVM.

\section{Métricas}

O projeto deve priorizar EER e min-DCF, não acurácia.

\subsection{EER}

EER é o ponto em que, de forma aproximada, taxas de erro de falso aceite e falsa rejeição se encontram. O valor deve ser calculado a partir dos scores e rótulos de teste, respeitando a partição.

\subsection{min-DCF}

Usar os parâmetros definidos no plano:
\[
C_{\mathrm{miss}}=C_{\mathrm{fa}}=1,\quad P_{\mathrm{tar}}=0.05.
\]

\section{Desempenho em CPU}

Registrar:
\begin{itemize}
\item tempo de leitura;
\item tempo de WPT;
\item tempo de reconstrução;
\item tempo de extração;
\item tempo de treino;
\item tempo de inferência;
\item memória.
\end{itemize}

O plano recomenda medir custo em CPU desde o início.

\section{Registro de experimentos}

Sugestão prática de tabela:

\begin{center}
\begin{tabular}{ll}
\toprule
Campo & Valor\\
\midrule
Data & \\
Corpus & \\
Wavelet & \\
Nível & \\
Modo de borda & \\
Classificador & \\
Seed & \\
Áudios & \\
Features & \\
Tempo & \\
Memória & \\
EER & \\
min-DCF & \\
Observações & \\
\bottomrule
\end{tabular}
\end{center}

\section{Estrutura recomendada do projeto}

\begin{lstlisting}
PIBITI/
|-- data/
|-- splits/
|-- src/
|-- features/
|-- results/
|-- reports/
|-- tests/
|-- notebooks/
|-- README.md
|-- requirements.txt
\-- .gitignore
\end{lstlisting}

Não colocar áudios dos corpora no repositório público quando as licenças/regras não permitirem sua redistribuição.

\section{Troubleshooting}

\begin{longtable}{p{3cm}p{4cm}p{7cm}}
\toprule
Sintoma & Possível causa & Diagnóstico/ação\\
\midrule
NaN nas features & Entrada inválida ou cálculo degenerado & Verificar áudio, intermediários e tratar casos explicitamente.\\
Reconstrução diferente & Padding/borda/alinhamento & Testar sinal sintético e registrar modo de borda.\\
Processamento lento & Corpus grande ou configuração pesada & Testar subconjunto, medir CPU e usar cache.\\
Memória insuficiente & Processamento em lote grande & Processar em lotes e salvar resultados intermediários.\\
Speaker leakage & Partição inadequada & Verificar locutor entre treino/validação/teste.\\
Resultados não reprodutíveis & Seed/configuração não registrada & Registrar seed, parâmetros e versão do ambiente.\\
\bottomrule
\end{longtable}

\section{Documentação semanal}

Sugestão de diário:

\begin{center}
\begin{tabular}{p{2cm}p{3cm}p{4cm}p{3cm}}
\toprule
Data & Atividade & Resultado & Próximo passo\\
\midrule
& & &\\
& & &\\
& & &\\
\bottomrule
\end{tabular}
\end{center}

\section{Relatório parcial}

Ao final do mês 6, o plano prevê relatório parcial e núcleo validado. Uma estrutura prática é:
\begin{enumerate}
\item Introdução;
\item Objetivos;
\item Referencial teórico;
\item Metodologia;
\item Corpora e partição;
\item WPT;
\item Descritores;
\item GMM e XGBoost;
\item Métricas;
\item Resultados;
\item Discussão;
\item Limitações;
\item Próximos passos.
\end{enumerate}

\section{Relatório final e apresentação}

O relatório final deve consolidar objetivos, metodologia, resultados, análise de sub-bandas, custo em CPU, limitações e trabalhos futuros. A apresentação deve mostrar problema, motivação, objetivo, metodologia, dados, WPT, features, classificadores, resultados e conclusão.

\section{Checklist geral}

\begin{itemize}[label=$\square$]
\item Ambiente reprodutível.
\item Repositório Git.
\item Bibliografia/fichamento.
\item Partição validada e congelada.
\item WPT db4 L=4 funcionando.
\item Reconstrução validada.
\item Features implementadas.
\item Matriz de features.
\item GMM.
\item XGBoost.
\item EER.
\item min-DCF.
\item Custo em CPU.
\item Expansão SVM/RF/db8/sym8/L=5.
\item Análise de sub-bandas.
\item Biblioteca documentada.
\item Artigo.
\item Relatório final.
\item Apresentação.
\end{itemize}

\section{O que fazer AGORA -- Mês 1, Semana 2}

\textbf{Hoje}
\begin{itemize}
\item Estudar o conceito de wavelet.
\item Entender wavelet-mãe.
\item Entender sub-bandas.
\item Entender DWT e WPT.
\end{itemize}

\textbf{Amanhã}
\begin{itemize}
\item Estudar árvore WPT.
\item Entender o nível $L$.
\item Confirmar por que $L=4$ produz 16 sub-bandas.
\item Fazer o pequeno teste com PyWavelets.
\end{itemize}

\textbf{Depois}
\begin{itemize}
\item Registrar o que aprendeu no fichamento.
\item Desenhar a árvore WPT.
\item Salvar o código de teste no repositório.
\item Registrar dúvidas para a reunião com o orientador.
\end{itemize}

\textbf{Ainda não fazer:}
\begin{itemize}
\item não treinar GMM;
\item não treinar XGBoost;
\item não calcular EER como etapa principal;
\item não processar o corpus inteiro apenas para testar a WPT;
\item não começar a expansão db8/sym8/L=5.
\end{itemize}

O critério prático para terminar a Semana 2 é conseguir explicar claramente o que é WPT, como funciona a árvore, o que significa $L$, por que $L=4$ gera 16 sub-bandas e como localizar essas sub-bandas com PyWavelets.

\section{Boas práticas permanentes}

O plano determina:
\begin{itemize}
\item validar a reconstrução antes de tudo;
\item manter uma partição única e congelada;
\item extrair features uma vez e reutilizar cache;
\item registrar nível, wavelet e borda;
\item medir custo em CPU;
\item versionar o trabalho em Git;
\item registrar corpus, nível, wavelet, classificador, seeds, métricas e tempos;
\item não misturar validação e teste;
\item respeitar split por locutor;
\item priorizar EER/min-DCF;
\item realizar reuniões quinzenais com o orientador;
\item escrever o relatório em paralelo ao trabalho.
\end{itemize}

\end{document}
